% MERA: Model Evolution and Routing with Skill Adaptation — AAAI 2027 submission
% paper.tex — auto-updated 2026-08-14 by weekly pipeline (v10)
% DO NOT edit the %% MANUAL-PIN blocks below.

\documentclass[letterpaper]{article}
\usepackage{aaai27}
\usepackage{times}
\usepackage{helvet}
\usepackage{courier}
\usepackage{url}
\usepackage{graphicx}
\usepackage{booktabs}
\usepackage{multirow}
\usepackage{amsmath}
\usepackage{amssymb}
\usepackage{microtype}
\usepackage{xcolor}
\usepackage{hyperref}
\usepackage[utf8]{inputenc}
\usepackage{amsthm}
\newtheorem{proposition}{Proposition}

\setlength{\pdfpagewidth}{8.5in}
\setlength{\pdfpageheight}{11in}

\title{MERA: Model Evolution and Routing with Skill Adaptation\\
for Agentic Systems at Scale}

\author{
Zeyu Wang\\
0G.ai / Institute of Artificial Intelligence\\
\texttt{zeyu.wang@0g.ai}
}

\begin{document}

\maketitle

%% ABSTRACT %%
\begin{abstract}
Large language model (LLM) serving faces a fundamental cost-quality trade-off: powerful
frontier models are expensive while cheaper models often fail on hard tasks.
We present \textbf{MERA} (Model Evolution and Routing with Skill Adaptation), a
self-improving agentic serving system that jointly evolves routing decisions and model
capabilities via reinforcement learning and supervised fine-tuning.
MERA maintains a \emph{SkillBook} of per-signature routing statistics and a
\emph{learned router} trained from execution traces.
On HumanEval (164 code tasks), MERA achieves \textbf{99\% task accuracy at 83\% lower
cost} than always routing to the frontier model.
A BERT-based router achieves \textbf{93.04\% routing accuracy} with a \textbf{2.12\%
fallback rate}.
Over 4 end-to-end evolution cycles on HumanEval with DAPO multi-turn repair ($G{=}8$),
the skills arm follows a \textbf{non-monotonic trajectory} (70.7\%$\to$65.9\%$\to$73.2\%$\to$75.6\%):
the cycle-1 dip is explained by SFT epoch-2 geometry conflict—the only cycle where
epoch-2 training loss \emph{increases} above epoch-1 (0.166$\to$0.271; Hypothesis F).
Mechanistic analysis reveals \textbf{52.4\% of training groups produce zero gradient} even with
DAPO dynamic sampling; without DAPO's $\sigma{=}0$ filter, these groups produce spurious
amplified gradients~\cite{darkroom2026}, not truly zero updates.
A standalone GRPO pass on MBPP yields +2pp pass@1 for Qwen2.5-Coder-1.5B ($n{=}1$ seed).
The Group-Standard-Deviation Identity~\cite{groupsd2026} provides formal grounding:
GRPO, Dr.~GRPO, and DAPO differ only in how they treat within-group $\sigma$; our
52.4\% zero-variance finding corresponds precisely to the $\sigma{=}0$ regime.
Extended to three-domain agentic tau2-bench tasks with a Qwen3.6-35B-A3B adapter,
MERA achieves \textbf{89.19\% task pass at 22.16\% of always-large cost} at peak;
a held-out evaluation shows the domain-specialized 35B model (80\%) surpassing the
frontier GPT-5.4 (71\%), revealing that agent specialization can render frontier
escalation counterproductive.
\end{abstract}

%% SECTION 1 %%
\section{Introduction}

Deploying LLMs at scale requires navigating a steep cost-quality frontier.
A single query to a frontier model costs 20--40$\times$ more than an equivalent query
to a small open-source model.
On diverse code-generation benchmarks, the small model solves $\approx$74\% of tasks
while the large model solves $\approx$95\%---leaving a stark deployment choice between
cost and quality.

The standard fallback strategy---always try the small model, escalate on failure---achieves
98\% accuracy at 73\% of the large-model cost.
But many tasks \emph{predictably} require the large model (complex algorithms, cryptographic
patterns), while others are reliably solved by the small model.
A \emph{routing} system that predicts this in advance avoids unnecessary large-model calls.

Beyond routing, \textbf{model evolution}---progressively training the small model on tasks
it currently fails---can expand the set of tasks the small model handles independently,
compounding cost savings over time.
Execution-grounded RL is a natural fit: the environment provides binary pass/fail rewards
with no human annotation.

We contribute:
\begin{enumerate}
\item \textbf{SkillBook routing evolution} (\S\ref{sec:skillbook}): online learning of
  per-signature routing statistics with Laplace-smoothed confidence estimates.
\item \textbf{Learned router} (\S\ref{sec:router}): BERT-based binary classifier achieving
  93.04\% routing accuracy vs.\ 68.28\% baseline.
\item \textbf{GRPO-based LLM evolution} (\S\ref{sec:grpo}): +2pp absolute MBPP pass@1 for
  Qwen2.5-Coder-1.5B.
\item \textbf{End-to-end ablation} (\S\ref{sec:ablation}): unified evaluation isolating
  each component's contribution.
\item \textbf{Multi-cycle MERA evolution} (\S\ref{sec:he4cyc}): 4-cycle DAPO joint
  evolution on HumanEval showing non-monotonic skills arm (70.7\%$\to$65.9\%$\to$73.2\%$\to$75.6\%)
  with cycle-1 dip explained by offline-confirmed SFT epoch-2 geometry conflict (Hypothesis~F)
  and 52.4\% zero-variance groups as the primary barrier to further gain.
\item \textbf{Failure mode analysis} (\S\ref{sec:analysis}): identification of
  winner-takes-all collapse, zero-variance silent groups, and SFT epoch-2 overtraining as
  the binding constraints; the Dark Room effect~\cite{darkroom2026} reveals that without
  DAPO's filter, zero-variance groups produce spurious amplified gradients rather than true zeros.
\item \textbf{Agentic extension} (\S\ref{sec:tau2}): 4-cycle joint evolution on tau2-bench
  showing the domain-specialized small model surpasses the frontier model on held-out
  tasks (80\% vs.\ 71\%), and identifying router overfitting in agentic settings.
\end{enumerate}

Tracks A (HumanEval skills routing), B (MBPP GRPO), C (tau2 joint evolution), and D
(4-cycle HumanEval DAPO) are ablation studies of MERA components at different
scales and modalities---not simultaneous deployments of a single system.

%% SECTION 2 %%
\section{Related Work}

\paragraph{LLM Routing.}
FrugalGPT~\cite{chen2023frugalgpt} routes queries to cheaper models via a learned cascade.
RouterBench~\cite{hu2024routerbench} benchmarks routing strategies across 11 LLMs.
Recent work explores regret-minimization routing~\cite{causal2026routing} and
skill-composable routing for agents~\cite{skillrouter2026}.
Our SkillBook extends routing with \emph{online} statistics from deployment traces,
requiring no offline re-training.

\paragraph{RL for Code Generation.}
DeepSeek-R1~\cite{deepseekr1} and GRPO~\cite{shao2024deepseekmath} establish that
verifiable binary rewards enable strong code-generation improvement.
We organize the rapidly growing literature by the failure mode each method addresses.

\textbf{(a) Zero-variance collapse and remedies.}
When all $K$ rollouts share the same reward, GRPO's $1/\sigma$ normalization amplifies
finite-precision noise to $\mathcal{O}(1)$ spurious gradient~\cite{darkroom2026}
(the Dark Room effect); DAPO~\cite{yu2025dapo} prevents this via an explicit $\sigma{=}0$
group filter.
REINFORCE++~\cite{zeng2025reinforceplusplus} uses a global EMA baseline providing
non-zero gradient on all-same-reward groups.
GCPO~\cite{chen2026gcpo} applies team-level AST-diversity coverage credit against
winner-takes-all collapse; MicroCoder~\cite{microcoder2026} uses diversity-determined
temperature; RECRL~\cite{recrl2026} applies adaptive curriculum to Goldilocks-zone tasks.
SRPO~\cite{srpo2026} routes all-fail rollouts to a self-distillation correction path,
converting zero-gradient all-fail groups to logit-level signal ($+3.2$pp HumanEval,
$+4.7$pp MBPP$+$ over DAPO).
Sign-advantage~\cite{signadv2026} replaces normalized advantage with $A{=}2r{-}1$,
guaranteeing non-zero gradient; its measured degeneracy rate ${\geq}0.69$ at $G{=}4$
is consistent with our 52.4\% ($G{=}8$, DAPO) finding.
The Group-Standard-Deviation Identity~\cite{groupsd2026} unifies GRPO, Dr.~GRPO, and
DAPO as three operations on $\sigma$: GRPO divides by $\sigma$ (unstable at
$\sigma{\to}0$), Dr.~GRPO drops the division (update $\propto p(1{-}p)$, naturally
weighted by informativeness), and DAPO filters $\sigma{=}0$ entirely.
We measure zero-variance rates empirically: 73\% at $G{=}4$ (MBPP), 52.4\% at $G{=}8$
DAPO (HumanEval).

\textbf{(b) Training stability and forgetting.}
Plasticity loss~\cite{plasticity2026} causes RL degradation in long training horizons
(empirically: 400-task GRPO degrades while 200-task improves).
SFT overtraining produces low-entropy checkpoints that cause gradient vanishing at
RL onset~\cite{sftovertraining2026}, directly explaining our cycle-1 skills arm dip
(Hypothesis~F, \S\ref{sec:analysis}).
Geometry conflict in LoRA subspace~\cite{geometry2026} between a mode-collapsed adapter
and subsequent SFT gradients causes epoch-2 loss inversion---confirmed in our training logs.
CPO~\cite{cpo2026} shows sequential RFT forgets diverse prior tasks; \citet{rftforgetting2026}
formalize that RFT regularization scales with within-group variance, so DAPO-zeroed groups
contribute neither learning \emph{nor} forgetting, explaining tau2 cycle oscillation.
DemoPSD~\cite{demopsd2026} modulates SFT loss by teacher-student disagreement rate.

\textbf{(c) Token-level credit assignment and data selection.}
EGCA~\cite{egca2026} assigns GRPO advantage to causally responsible token spans via
execution-trace divergence ($+1.5$pp MBPP).
TAROT~\cite{tarot2026} introduces intra-problem test-tier curriculum for code RL.
GMPO~\cite{gmpo2026} replaces arithmetic-mean token aggregation with geometric mean,
bounding gradient distortion at syntactically critical positions ($+4.1\%$ math).
Reward-guided dataset distillation~\cite{zpddistill2026} filters SFT traces to the
zone of proximal development ($r_\text{small}\in[0.2,\,0.8]$),
yielding $+3.1$--$5.6$pp over naive SFT; we apply this analysis in
\S\ref{sec:analysis} to explain our 20.7pp skills gap.

%% SECTION 3 %%
\section{SkillBook: Signature-Based Routing Evolution}
\label{sec:skillbook}

\subsection{Prompt Signatures}

We extract a compact signature from each incoming prompt:
$\texttt{sig} = \langle\texttt{tier}\rangle\,|\,\langle\texttt{tags}\rangle$,
where $\texttt{tier}\in\{S,M,L\}$ encodes prompt length ($<$200, 200--500, $>$500 chars)
and $\texttt{tags}$ is a set extracted by keyword matching (e.g., \texttt{list},
\texttt{num}, \texttt{sort}, \texttt{crypto}).

\subsection{SkillBook Learning}

For each (signature, model) pair the SkillBook maintains success and total counts.
Routing confidence uses Laplace smoothing:
\begin{equation}
  \hat{p}(\texttt{sig},m) = \frac{n^+_{\texttt{sig},m}+1}{n_{\texttt{sig},m}+2}
\end{equation}
New entries start at 0.5 (agnostic).
The SkillBook recommends the cheapest model whose confidence exceeds a target threshold.

\subsection{Escalation Policy}

The policy routes to the recommended model, executes generated code against unit tests,
escalates to the large model on failure, and records outcomes back into the SkillBook.
This guarantees accuracy $\geq$ always-large while reducing cost.

\subsection{Results}

\begin{table}[h]
\centering
\caption{HumanEval 164-task strategy comparison.}
\label{tab:humaneval}
\begin{tabular}{lrrr}
\toprule
Strategy & Accuracy & Cost & vs.\ Large \\
\midrule
Always-large (GPT-5.4) & 95\% & \$0.37 & 100\% \\
Always-small (DeepSeek) & 88\% & \$0.01 & 4\% \\
Router only & 74\% & \$0.15 & 42\% \\
Router + Fallback & 98\% & \$0.27 & 73\% \\
\textbf{Router + Skills Evolve} & \textbf{99\%} & \textbf{\$0.06} & \textbf{17\%} \\
\bottomrule
\end{tabular}
\end{table}

Router + Skills Evolve achieves the highest accuracy (99\%) at the lowest cost (17\% of
always-large). Across 4 evolution rounds, the SkillBook learns 49 signatures: 40 downgrade
to the small model, 9 are flagged as large-model-required.

%% SECTION 4 %%
\section{Learned Router}
\label{sec:router}

\subsection{Training Data}

Routing labels are derived from execution traces: small-model pass $\to$ label \texttt{small};
small-model fail + large-model pass $\to$ label \texttt{large}.
We augment with 3{,}328 UncommonRoute bench examples.
Evolver Dataset v1 statistics: 2{,}662 train / 332 dev / 334 test router examples;
2{,}258 small-route / 1{,}070 large-route labels.

\subsection{Model and Threshold Tuning}

We fine-tune BERT-base-uncased (110M parameters) for 8 epochs ($\eta=2\times10^{-5}$).
The threshold $\theta$ is tuned to satisfy a $\leq 2\%$ fallback constraint on the dev split,
yielding $\theta^*=0.57$.
BERT inference adds ${<}10$ms/query at batch size~1 on CPU, adding negligible serving latency.

\subsection{Results}

\begin{table}[h]
\centering
\caption{End-to-end routing ablation (848 evaluation examples).}
\label{tab:routing}
\begin{tabular}{lrrrr}
\toprule
Variant & Acc. & Large F1 & Fallback & Cost \\
\midrule
Base (always-small) & 68.28\% & 0.00\% & 31.72\% & 33.54\% \\
+~SkillBook & 69.46\% & 24.93\% & 26.65\% & 37.27\% \\
\textbf{+~Learned router} & \textbf{93.04\%} & \textbf{89.48\%} & \textbf{2.12\%} & 37.75\% \\
\bottomrule
\end{tabular}
\end{table}

The learned router reduces fallback from 31.72\% to 2.12\% and raises accuracy by 24.76pp.
SkillBook alone contributes modestly (+1.18pp accuracy).

%% SECTION 5 %%
\section{LLM Evolution via GRPO}
\label{sec:grpo}

\subsection{Setup}

We train Qwen2.5-Coder-1.5B-Instruct with local GRPO: 200 MBPP augmented train tasks,
4 rollouts per prompt, 1 epoch, LoRA rank 16, $\eta=5\times10^{-6}$, max 192 new tokens.
Reward is binary (executable test pass/fail). We use qwen-chat prompt format.
We use no KL penalty ($\beta{=}0$) and clip ratios $\varepsilon_\text{low}{=}0.2$,
$\varepsilon_\text{high}{=}0.5$ (DAPO configuration).

\subsection{Scaling Results}

\begin{table}[h]
\centering
\caption{LLM evolution: scaling attempts on MBPP eval100.}
\label{tab:llm}
\begin{tabular}{llrr}
\toprule
Model & Method & Base & Trained \\
\midrule
Qwen2.5-Coder-0.5B & SFT 120 tasks & 10/20 & 9/20 \\
Qwen2.5-Coder-0.5B & GRPO 40$\times$4 & 8/20 & 9/20 \\
\textbf{Qwen2.5-Coder-1.5B} & \textbf{GRPO 200$\times$4} & \textbf{47/100} & \textbf{49/100} \\
Qwen2.5-Coder-1.5B & GRPO 400$\times$4 & 47/100 & 46/100 \\
Qwen2.5-Coder-3B & GRPO 200$\times$4 & 62/100 & 60/100 \\
Qwen2.5-Coder-3B & SFT 400 tasks & 62/100 & 62/100 \\
Qwen2.5-Coder-7B & GRPO 100$\times$4 & 77/100 & 75/100 \\
\bottomrule
\end{tabular}
\end{table}

Naive scaling of training budget or model size consistently fails.
The best positive result is 1.5B GRPO 200$\times$4 (+2pp, $n{=}1$ seed).

\subsection{Multi-Cycle HumanEval Evolution (DAPO)}
\label{sec:he4cyc}

We run a 4-cycle end-to-end MERA pipeline on HumanEval (82 tasks) using GPT-5.5 as the
large model and a Qwen3.6-35B-A3B adapter as the evolving small model, with DAPO
multi-turn repair rollout ($G{=}8$, max\_turns$=3$, LoRA $r{=}16$, $\eta{=}5\times10^{-6}$).

\begin{table}[h]
\centering
\caption{4-cycle HumanEval MERA---per-cycle results ($n{=}1$ seed;
  $^\dagger$binomial SE ${\approx}{\pm}5.4$pp at 82 tasks).}
\label{tab:he4cyc}
\begin{tabular}{lrrrr}
\toprule
Cycle & Full Pass$^\dagger$ & Route Acc$^\dagger$ & Cost & Skills Arm$^\dagger$ \\
\midrule
0 & 91.46\% & 92.68\% & 31.95\% & 70.73\% \\
1 & 92.68\% & 90.24\% & 40.73\% & 65.85\%$^*$ \\
2 & 91.46\% & 93.90\% & 28.66\% & 73.17\% \\
3 & 92.68\% & 92.68\% & 27.56\% & \textbf{75.61\%} \\
\midrule
Always-large (GPT-5.5) & \multicolumn{2}{r}{96.34\%} & 100\% & --- \\
\bottomrule
\end{tabular}
{\small $^*$Cycle-1 SFT overfit: epoch-2 loss $0.166{\to}0.271$; epoch-1 ckpt saves 4.88pp.}
\end{table}

The skills arm follows a non-monotonic trajectory: 70.73\%$\to$65.85\%$\to$73.17\%$\to$\textbf{75.61\%},
improving 4.88pp overall yet dipping at cycle~1.
The cycle-1 dip is anti-correlated with ACR (47.6\%—the lowest ACR of all 4 cycles),
ruling out zero-variance collapse as the sole cause and suggesting SFT quality as the
primary bottleneck.
Offline analysis of cycle-1 SFT training logs confirms Hypothesis~F:
epoch-2 loss \emph{increases} from 0.166 to 0.271 (overfit\_gap\,$=$\,$+0.105$; the
only cycle with positive overfit\_gap), coinciding with entropy compression from 0.255
to 0.207---consistent with geometry conflict~\cite{geometry2026} between the mode-collapsed
GRPO-0 adapter and cycle-1 SFT data. Using the epoch-1 checkpoint (checkpoint-6,
available in the repo) as the GRPO initializer instead of the overtrained epoch-2 checkpoint
is predicted to recover the 4.88pp dip (EXP-137).
The Full system equals the Router-only system in all 4 cycles (both 92.68\% at cycle~3):
routing is a prompt-level decision; GRPO improves code accuracy in the non-collapsed
47.6\% of groups but does not shift the set of prompts requiring large-model escalation.

\paragraph{Zero-Variance Diagnosis (Cycle~3 GRPO).}
Inspection of the cycle-3 GRPO training log reveals a concrete bottleneck: of 82
training groups ($G{=}8$, DAPO dynamic sampling), \textbf{43/82 (52.4\%) were dropped
as zero-variance}---despite DAPO's post-hoc filtering. The informative subset:
39 groups, 312 rollouts used for weight updates. Zero-variance breaks down as
${\approx}40\%$ all-pass (problems the cycle-3 model already masters) and
${\approx}12\%$ all-fail (intractable at 1.5B capacity). This is lower than the
73\% zero-variance at $G{=}4$ standard GRPO (§\ref{sec:analysis}): DAPO reduces
but does not eliminate the bottleneck.

\begin{proposition}[Zero-Gradient Groups under DAPO]
\label{prop:zero-grad}
Under DAPO with dynamic sampling (drop groups where $\sigma{=}0$), let $K$ denote
rollout count per group. For any group where all $K$ outcomes are identical
(all-pass or all-fail), within-group reward variance $\sigma^2{=}0$ and DAPO drops the
group. By the Group-SD Identity~\cite{groupsd2026}, the DAPO policy gradient for that
group is exactly zero. Without DAPO's filter, standard GRPO's $1/\sigma$ normalization
amplifies finite-precision noise ($\sigma{\approx}\varepsilon$) to $\mathcal{O}(1)$
spurious gradient~\cite{darkroom2026}. In our 4-cycle HumanEval experiments, DAPO
correctly zeros 46.3\%--52.4\% of training groups per cycle; the remaining 47.6\%--53.7\%
carry all informative gradient.
\end{proposition}

%% SECTION 6 %%
\section{End-to-End Ablation}
\label{sec:ablation}

\begin{table}[h]
\centering
\caption{Component ablation. Routing on 848 examples; code pass on MBPP eval100.}
\label{tab:ablation}
\begin{tabular}{lcccrrrr}
\toprule
Variant & Sk. & Rt. & LL & Acc. & FB & Cost & Pass \\
\midrule
Base & -- & -- & -- & 68.3\% & 31.7\% & 33.5\% & 47\% \\
+Skills & \checkmark & -- & -- & 69.5\% & 26.7\% & 37.3\% & 47\% \\
+Router & \checkmark & \checkmark & -- & 93.0\% & 2.1\% & 37.8\% & 47\% \\
Full & \checkmark & \checkmark & \checkmark & \textbf{93.0\%} & \textbf{2.1\%} & 37.8\% & \textbf{49\%} \\
\bottomrule
\end{tabular}
\end{table}

Router training is the single largest contributor; LLM GRPO adds orthogonal code quality gain
(+2pp MBPP pass@1).
The Full=Router routing result holds because routing quality is a \emph{prompt-level} decision:
the GRPO adapter improves code accuracy on the tasks it handles but does not shift the
distribution of prompts requiring large-model routing.

%% SECTION 6.5 — AGENTIC EXTENSION %%
\section{Agentic Extension: Tau2-Bench Joint Evolution}
\label{sec:tau2}

We extend MERA to tau2-bench, a three-domain agentic customer-service benchmark
(airline, retail, telecom). The ``small'' model is a Qwen3.6-35B-A3B adapter
fine-tuned via SFT on hard agent traces; the ``large'' model is GPT-5.4.
We run the full MERA pipeline (SkillBook $\to$ LLM SFT $\to$ Router) for 4 cycles.

\subsection{Joint Evolution Run (74 In-Distribution Tasks)}

\begin{table}[h]
\centering
\caption{Tau2 joint evolution per-cycle results (74 in-sample tasks,
Qwen3.6-35B-A3B adapter).}
\label{tab:tau2_joint}
\begin{tabular}{lrrrr}
\toprule
Cycle & Task Pass & Router Acc & Fallback & Cost vs.\ Large \\
\midrule
Cycle 0 & 78.38\% & 93.24\% & 4.05\% & 47.70\% \\
\textbf{Cycle 1} & \textbf{89.19\%} & 93.24\% & \textbf{4.05\%} & \textbf{22.16\%} \\
Cycle 2 & 70.27\% & 82.43\% & 5.41\% & 48.92\% \\
Cycle 3 & 72.97\% & 89.19\% & 6.76\% & 35.54\% \\
\bottomrule
\end{tabular}
\end{table}

At peak (Cycle~1), MERA achieves 89.19\% task pass at 22.16\% of always-large cost---a
4.5$\times$ cost reduction with competitive accuracy.

\subsection{Formal Train/Test Split (178 Train / 100 Held-Out)}

We run a second experiment with zero train/test overlap (all three domains: retail 40,
telecom 40, airline 20 held-out tasks).

\begin{table}[h]
\centering
\caption{Tau2 held-out evaluation (100 unseen tasks, final-cycle adapter).}
\label{tab:tau2_heldout}
\begin{tabular}{lrrr}
\toprule
Variant & Task Pass & Route Acc & Cost vs.\ Large \\
\midrule
Always-large (GPT-5.4) & 71.00\% & --- & 100\% \\
Base (35B MoE, always-small) & \textbf{80.00\%} & 80.00\% & \textbf{10\%} \\
+~SkillBook & 80.00\% & 80.00\% & 10\% \\
+~Router & 71.00\% & 66.00\% & 46\% \\
\bottomrule
\end{tabular}
\end{table}

\paragraph{Tau2 cycle instability and SFT forgetting.}
The joint-run oscillation (89.19\%$\to$70.27\%$\to$72.97\%) is consistent with
SFT-driven catastrophic forgetting~\cite{rftforgetting2026}: SFT lacks RFT's implicit
variance-proportional regularization, making sequential SFT over successive cycle traces
prone to forgetting earlier cycle improvements. Converting the tau2 adapter training from
SFT to GRPO/RFT is queued as EXP-111.

\paragraph{Key findings.}
\textbf{(1) Specialized small model outperforms frontier.}
On held-out tasks, the domain-fine-tuned 35B MoE achieves 80\% task pass, exceeding
GPT-5.4 at 71\%, at only 10\% of the cost.
\textbf{(2) Router overfits in agentic domains.}
The TF-IDF+LR router drops to 66\% routing accuracy on held-out, 14pp below the
always-small baseline (80\%). Routing to the large model actively \emph{hurts}:
task pass degrades from 80\% to 71\% while cost rises from 10\% to 46\%.
This is a fundamental limitation of in-distribution routing on agentic multi-turn tasks
where frontier model quality is not reliably superior.
\textbf{(3) SkillBook provides no gain beyond base.}
With the 35B adapter already specialized, the signature-based SkillBook cannot improve
beyond always-small routing.

%% SECTION 7 %%
\section{Failure Mode Analysis}
\label{sec:analysis}

\paragraph{Zero-Variance Silent-Group Problem and Dark Room Effect.}
With $G{=}4$ rollouts and binary reward on MBPP, 53\% of tasks yield all-fail groups
and 20\% yield all-pass groups. Standard GRPO advantage normalization
$A_i=(r_i-\bar{r})/\sigma_r$ produces $0/0$ on all-same-reward groups, nominally
clipped to zero. However, the Dark Room effect~\cite{darkroom2026} reveals a critical
pathology: when $\sigma\to 0$ without an explicit filter, finite-precision residuals
($\sigma{\approx}\varepsilon$) are amplified to $\mathcal{O}(1)$ spurious gradients that
actively \emph{degrade} performance rather than contributing zero.
DAPO's explicit $\sigma{=}0$ drop (Proposition~\ref{prop:zero-grad}) prevents this:
it converts spurious amplified updates to true zeros.
Scaling to $G{=}8$ with DAPO on HumanEval reduces the zero-variance fraction to
\textbf{52.4\%} (43/82 groups; §\ref{sec:he4cyc}), leaving 39 informative groups.
Two complementary analyses confirm the group-level picture.
\citet{rftforgetting2026} show RFT's implicit regularization scales with within-group
variance; DAPO-zeroed groups contribute neither learning nor forgetting.
The Group-SD Identity~\cite{groupsd2026} unifies GRPO/Dr.~GRPO/DAPO as three operations
on $\sigma$: Dr.~GRPO (no $1/\sigma$) provides update $\propto p(1{-}p)$---naturally
zero for collapsed groups and maximal at $p{=}0.5$.

\paragraph{SFT Epoch-2 Geometry Conflict (Hypothesis~F).}
Offline analysis of cycle-1 SFT training logs (available in the repo:
\texttt{results/e2e\_4cyc\_gpt55/cycle\_1/llm\_adapter/training\_info.json})
reveals a novel failure mode: epoch-2 loss increases from 0.166 to 0.271
(overfit\_gap\,$=$\,$+0.105$), the \emph{only} cycle with a positive overfit gap.
Cycle-1 SFT starts from the GRPO-0 adapter (entropy 0.152, mode-collapsed), whose
compressed parameter space conflicts with cycle-1 teacher-trace gradients in epoch~2
(geometry conflict~\cite{geometry2026}).
The epoch-1 checkpoint (checkpoint-6; $\ell_1{=}0.166$, entropy 0.255) is saved in
the repo and provides a better GRPO initializer.
This mechanistically explains the cycle-1 skills arm dip (65.85\% vs.\ 70.73\%):
the overtrained SFT checkpoint reduces model plasticity~\cite{sftovertraining2026},
causing gradient vanishing at the start of GRPO-1.
Selecting the epoch-1 checkpoint when overfit\_gap\,$>\,0$ is predicted to recover
the 4.88pp dip with a 1-line change to the pipeline (EXP-137).

\paragraph{Skills Gap.}
The skills arm (always-small $+$ procedure) achieves 75.61\% pass@1 on HumanEval,
a 20.7pp gap below always-large (96.34\%).
The bottleneck is SFT data quality: current training uses all ~77 teacher traces
including tasks the small model already solves ($r_\text{small}{>}0.8$) and tasks
beyond its reach ($r_\text{small}{<}0.2$).
Reward-guided dataset distillation~\cite{zpddistill2026} addresses this by filtering
traces to the zone of proximal development ($r_\text{small}\in[0.2, 0.8]$), reporting
$+3.1$--$5.6$pp over naive SFT on code benchmarks.
Applying ZPD filtering to our HumanEval pipeline would retain ${\approx}47$ of 77
traces (removing ${\approx}18$ trivially solved and ${\approx}12$ intractable tasks),
concentrating teacher supervision where the student can most benefit.
This is queued as EXP-121 (ZPD-SFT) and EXP-122 (Joint ZPD-SFT $+$ Dr.~GRPO).

\paragraph{Winner-Takes-All Collapse.}
On mixed-outcome tasks, GRPO reinforces the single passing rollout until it dominates
subsequent rollout batches. Group advantage collapses to zero; the model converges to a
single solution mode per task.

\paragraph{Remedies.}
REINFORCE++~\cite{zeng2025reinforceplusplus} uses a global EMA baseline providing
gradient on all-same-reward groups.
GCPO~\cite{chen2026gcpo} applies team-level AST-diversity coverage to prevent
winner-takes-all collapse.
RECRL~\cite{recrl2026} and MicroCoder~\cite{microcoder2026} reduce zero-variance via
adaptive curriculum and diversity-determined temperature.
SRPO~\cite{srpo2026} routes failed rollouts to a self-distillation correction path.
Sign-advantage~\cite{signadv2026} eliminates zero-gradient groups ($A{=}2r{-}1$).
Dr.~GRPO~\cite{groupsd2026} drops $1/\sigma$ entirely; update $\propto p(1{-}p)$.
A Dr.~GRPO variant is a 5-line change to \texttt{grpo\_train\_simple.py}; GPU access is
currently blocked (queued as EXP-120 for camera-ready validation).
No-std-norm~\cite{darkroom2026} (remove only the $\sigma$ division, keep mean subtraction)
converts spurious amplified gradients in near-zero-$\sigma$ groups to exact zeros,
the minimal fix to the Dark Room pathology and complementary to DAPO's group-level filter.
GMPO~\cite{gmpo2026} targets non-collapsed group quality via geometric-mean token aggregation.
For the skills gap, ZPD distillation~\cite{zpddistill2026} filters SFT traces to the
zone of proximal development; epoch-1 SFT checkpoint selection (Hypothesis~F, EXP-137)
provides an orthogonal and immediately testable fix.
These are queued as priority experiments pending A800 restoration.

%% SECTION 8 %%
\section{Conclusion}

We present MERA, achieving 99\% task accuracy at 17\% of frontier cost on HumanEval via
skills evolution, 93.04\% routing accuracy with a learned BERT router, and +2pp MBPP
pass@1 via GRPO (standalone; $n{=}1$ seed).
Over 4 end-to-end HumanEval evolution cycles, the skills arm follows a non-monotonic
trajectory (70.7\%$\to$65.9\%$\to$73.2\%$\to$75.6\%); offline analysis confirms the
cycle-1 dip is caused by SFT epoch-2 geometry conflict (Hypothesis~F), recoverable by
selecting the epoch-1 checkpoint---a 1-line fix testable without new GPU runs.
A mechanistic analysis pinpoints 52.4\% zero-variance DAPO groups as the primary
gradient barrier; Proposition~\ref{prop:zero-grad} formalizes the DAPO safety guarantee
and contrasts it with the Dark Room spurious-gradient pathology of standard GRPO.
Extended to tau2 agentic tasks, MERA reaches 89.19\% task pass at 22.16\% of always-large
cost at peak, and reveals that the domain-specialized small model (80\%) surpasses
the frontier GPT-5.4 (71\%) on held-out tasks, reframing routing as a domain-specialization
decision.
We identify SFT epoch-2 overtraining, router overfitting, and 52\%-zero-variance GRPO
as the four binding constraints (with winner-takes-all collapse), each with a concrete
low-overhead remedy grounded in both theory and our offline-confirmed experimental evidence.

%% REFERENCES %%
\bibliographystyle{aaai}
\bibliography{references}

\end{document}
